\section{Height two from two point counts}
\label{sec:point-count-height-two}
Let $X/\mathbb F_p$ be a smooth geometrically integral hypersurface
of degree $d\ge4$ in $\mathbb P^{d-1}$, and put $n=d-2$.
Its top Artin--Mazur formal group is one-dimensional.
Write $N_r=\#X(\mathbb F_{p^r})$ and
\begin{equation}
 \tau_r=(-1)^n\left(N_r-\sum_{i=0}^{n}p^{ir}\right).
                                                               \label{eq:PC1}
\end{equation}
The Frobenius convention is the one in the point-count trace formula.

\subsection{Classical inputs with their hypotheses}
Weak Lefschetz and duality identify all cohomology outside the middle
with the corresponding projective-space classes. One can recover this
particular weak Lefschetz statement from affine vanishing
\citep[Theorem~15.1, p.~105]{MilneLEC2013} and Poincar\'e duality
\citep[Theorem~24.1, p.~144]{MilneLEC2013}.
Indeed, the complement $U=\mathbb P^{n+1}\setminus X$ is smooth affine
of dimension $n+1$. Duality and affine vanishing give
$H_c^j(U,\mathbb Q_\ell)=0$ for $j<n+1$.
The compact-support exact sequence therefore makes
$H^i(\mathbb P^{n+1})\to H^i(X)$ an isomorphism for $i<n$ and
an injection for $i=n$. Duality on $X$ treats the degrees above $n$.
When $n$ is even, the middle ambient class has nonzero self-pairing
equal to the degree, so its orthogonal complement is Frobenius stable.
When $n$ is odd there is no middle ambient class.

Consequently, \eqref{eq:PC1} is the trace of the $r$th power of Frobenius on
primitive middle cohomology. This exact complete-intersection
decomposition and trace computation also appear in the proof of
Deligne's Theorem~(8.1), pp.~301--302; the trace formula is
(1.5.1), p.~275 \citep{Deligne1974}.
Deligne's Theorem~(1.6), p.~276, makes the middle characteristic
polynomial integral. Dividing out the middle ambient factor, when
present, preserves integrality.

Katz--Messing's Theorem~1, p.~74 \citep{KatzMessing1974}, identifies
the characteristic polynomials in crystalline and $\ell$-adic
cohomology for smooth projective absolutely irreducible varieties.
The crystalline theory is explicitly included in the statement's
preceding paragraph. Thus the $p$-adic valuations of the primitive
reciprocal roots are its crystalline slopes.

For these smooth complete intersections, the Artin--Mazur comparison
hypotheses were verified above. III, Corollary~3.4, p.~119
\citep{AM}, describes the slopes $a_i/b_i<1$, written
with multiplicity $b_i$, by
\[
 \sum_i(b_i-a_i)\le h^{0,n}=1,
\]
with equality exactly in finite formal height.
Each nonempty summand is a positive integer. In finite height there
is therefore just one slope below one; since its multiplicity is the
height, it is $(h-1)/h$ with multiplicity $h$.
In infinite height the inequality is strict, so there is no slope
below one. The primitive part has the same slopes below one: the
discarded middle ambient class, if any, has slope $n/2\ge1$.

\begin{theorem}[Point-count test on the Hasse-zero locus]
\label{thm:height-two-point-counts}
Assume that the formal height of $X$ is at least two.
Set
\[
 e_2=\frac{\tau_1^2-\tau_2}{2}\in\mathbb Z.
\]
Then
\begin{equation}
 h(X)=2\quad\Longleftrightarrow\quad v_p(e_2)=1,       \label{eq:PC2}
\end{equation}
where $v_p(0)=+\infty$.
For odd $p$, this is equivalent to
\begin{equation}
 \tau_2/p\not\equiv0\pmod p,                         \label{eq:PC3}
\end{equation}
which requires $N_2$ only modulo $p^2$, since
$\sum_{i=0}^{n}p^{2i}\equiv1\pmod{p^2}$.
For $p=2$, the equivalent test is
\begin{equation}
 \frac{\tau_1^2-\tau_2}{4}\not\equiv0\pmod2,          \label{eq:PC4}
\end{equation}
which is determined by $N_1\bmod4$ and $N_2\bmod8$.
\end{theorem}
\begin{proof}
Let $\alpha_1,\ldots,\alpha_B$ be the primitive Frobenius eigenvalues
in an algebraic closure of $\mathbb Q_p$, counted with multiplicity.
Newton's identity gives
\[
 e_2=\sum_{i<j}\alpha_i\alpha_j
     =(\tau_1^2-\tau_2)/2,
\]
and the integral characteristic polynomial proves $e_2\in\mathbb Z$.

If $h=2$, precisely two eigenvalues have valuation $1/2$, while all
remaining eigenvalues have valuation at least one. Their product is
the unique term of valuation one in $e_2$; every other term has
valuation strictly greater than one. Therefore $v_p(e_2)=1$.

If $3\le h<\infty$, every eigenvalue has valuation at least
$1-1/h>1/2$, so every pair product has valuation greater than one.
Integrality implies $v_p(e_2)\ge2$. If $h=\infty$, every valuation is
at least one, which gives the same conclusion.
This proves \eqref{eq:PC2}.

Under the hypothesis $h\ge2$, every eigenvalue has positive valuation.
Therefore the integer $\tau_1$ is divisible by $p$.
For odd $p$,
\[
 e_2\equiv-\tau_2/2\pmod{p^2},
\]
so \eqref{eq:PC2} is equivalent to \eqref{eq:PC3}. The quotient is integral, since
$\tau_2=\tau_1^2-2e_2$ is divisible by $p$.

At $p=2$, $e_2$ is even under the same hypothesis, and
$e_2/2=(\tau_1^2-\tau_2)/4$.
An even $\tau_1$ has square modulo eight determined by its residue
modulo four. Thus \eqref{eq:PC4} uses exactly the stated sufficient precisions.
In terms of counts alone, because $n\ge1$,
\begin{equation}
 \frac{(N_1-3)^2-(-1)^n(N_2-5)}4\pmod2             \label{eq:PC5}
\end{equation}
is the same test: the ambient sums are $3\bmod4$ and $5\bmod8$.
\end{proof}

The height-at-least-two hypothesis is essential to this formulation.
The characteristic-two formula cannot in general be reduced to
point counts modulo $p^2=4$, as the following smooth examples show.

\subsection{An exact characteristic-two precision obstruction}
Over $\mathbb F_2$, in variables $x_1,x_2,x_3,y,z$, define
\[
 F_A=yzA+y^5+x_1^5+x_2^5+x_3^5-x_1z^4,
\]
with either
\[
 A_0=x_1^3,\qquad A_2=x_1^3+x_2^3+x_3^3.
\]
Both equations define smooth quintic threefolds.
Indeed,
\[
 (F_A)_y=zA+y^4,\qquad (F_A)_z=yA,\qquad
 y(F_A)_y-z(F_A)_z=y^5.
\]
At a common zero of all partial derivatives, $y=0$ and the other
partials force $x_2=x_3=0$ and $x_1^4=z^4$.
At a projective point this makes $x_1,z$ nonzero, whereas
$(F_A)_y=z x_1^3\ne0$, a contradiction.

For both forms $a_2=0$. For $a_4$, whose target has exponent three
in every variable, none of the pure fifth powers or the term
$x_1z^4$ can occur. Consequently
\[
 a_4(F_A)=[(x_1x_2x_3)^3]A^3
 =\begin{cases}0,&A=A_0,\\6,&A=A_2.\end{cases}
\]
The logarithm criterion therefore gives height at least three for
$X_{A_0}$ and height exactly two for $X_{A_2}$.

Exact enumeration over $\mathbb F_2$ and
$\mathbb F_4=\mathbb F_2[t]/(t^2+t+1)$ gives
\[
\begin{array}{c|rrrrrr}
 A&N_1&N_2&\tau_1&\tau_2&e_2&h\\\hline
 A_0&15&85&0&0&0&\ge3\\
 A_2&15&89&0&-4&2&2
\end{array}
\]
Projective representatives were enumerated over both fields, and the
identity $C_q=1+(q-1)N_q$ was checked separately on the affine cones.
The reproducible integer and finite-field calculations are included
with the other programs in the workbench.


The examples have the same $N_1$ even as an integer and the same
$N_2\bmod4$, but different height-two status. Hence those coarser
point-count data do not decide height two, even within smooth
anticanonical hypersurfaces. Their $N_2\bmod8$ values differ, as
the proved criterion requires.
